\item[\underline{\underline{Transaction-row n$^°~(\bm{i + \roffTxInitTxn})$:}}]
	we must justify certain predictions made by the \userTxnDataMod{} module:
	\begin{description}
		\item[\underline{(Partially) justifying \txRequiresEvmExecution{}:}]
			we impose the following
			\begin{enumerate}
				\item $\txRequiresEvmExecution_{i + \roffTxInitTxn} = \rOne$;
				\item \If $\locTxInitIsMessageCall = 1$ \Then $\locNonSkipMessageCall = 1$ % $\accHasCode     _{i + \roffTxInitAccDoingRecipientValueReception}  =    1$
				\item \If $\locTxInitIsDeployment  = 1$ \Then $\txInitCodeSize _{i + \roffTxInitTxn}       \neq 0$ (\trash)
			\end{enumerate}
			where we set
			\[
				\left\{ \begin{array}{lcl}
					\locNonSkipMessageCall & \define & 
					\left[ \begin{array}{cl}
						+ & \locRecipientHasCodeAndIsntDelegated                       \\
						+ & \locRecipientIsDelegatedAndDelegateHasCodeAndIsntDelegated \\
					\end{array} \right]
					\vspace{2mm} \\
					\locRecipientHasCodeAndIsntDelegated                       & \define &
					\left[ \begin{array}{cl}
						\cdot & \locRecipientHasNonemptyCode \\
						\cdot & \locRecipientIsntDelegated   \\
					\end{array} \right]
					\vspace{2mm} \\
					\locRecipientIsDelegatedAndDelegateHasCodeAndIsntDelegated & \define &
					\left[ \begin{array}{cl}
						% \cdot & \locRecipientHasNonemptyCode \\
						\cdot & \locRecipientIsDelegated     \\
						\cdot & \locDelegateHasNonemptyCode  \\
						\cdot & \locDelegateIsntDelegated    \\
					\end{array} \right]
				\end{array} \right.
			\]
			where we have used the following auxiliary shorthands
			\[
				\left\{ \begin{array}{lcl}
					\locSenderHasNonemptyCode    & \define & \accHasCode     _{i + \roffTxInitAccDoingSenderPayForGas         } \\
					\locRecipientHasNonemptyCode & \define & \accHasCode     _{i + \roffTxInitAccDoingRecipientValueReception } \\
					\locDelegateHasNonemptyCode  & \define & \accHasCode     _{i + \roffTxInitAccViewingDelegationOrRecipient } \\
					\locSenderIsDelegated        & \define & \accIsDelegated _{i + \roffTxInitAccDoingSenderPayForGas         } \\
					\locRecipientIsDelegated     & \define & \accIsDelegated _{i + \roffTxInitAccDoingRecipientValueReception } \\
					\locDelegateIsDelegated      & \define & \accIsDelegated _{i + \roffTxInitAccViewingDelegationOrRecipient } \\
				\end{array} \right.
			\]
			as well as their ``negations''
			\[
				\left\{ \begin{array}{lcl}
					\locSenderHasEmptyCode     & \define & 1 - \locSenderHasNonemptyCode    \\
					\locRecipientHasEmptyCode  & \define & 1 - \locRecipientHasNonemptyCode \\
					\locDelegateHasEmptyCode   & \define & 1 - \locDelegateHasNonemptyCode  \\
					\locSenderIsntDelegated    & \define & 1 - \locSenderIsDelegated        \\
					\locRecipientIsntDelegated & \define & 1 - \locRecipientIsDelegated     \\
					\locDelegateIsntDelegated  & \define & 1 - \locDelegateIsDelegated      \\
				\end{array} \right.
			\]
			\saNote{}
			We remind the reader that delegated accounts ($\accIsDelegated \equiv \true$)
			automatically have nonempty code ($\accHasCode \equiv \true$), see
			constraint~(\ref{hub: account: account delegation: delegated accounts have code and code size equal to 23}) and
			constraint~(\ref{hub: account: account delegation: newly delegated accounts have code and code size equal to 23}).
		\item[\underline{Justifying $\txFinalRefundCounter$:}]
			cannot be set at the moment;
		\item[\underline{Justifying $\txInitialBalance$:}]
			we impose that $\txInitialBalance _{i + \roffTxInitTxn} = \accBalance _{i + \roffTxInitAccDoingSenderPayForGas}$
		\item[\underline{Justifying \txStatusCode{}:}]
			cannot be set at the moment\footnote{We \textbf{could} actually set the status code of the transaction at the present time as it coincides with  \locTransactionWontRevert{}.};
			% \item[\underline{Justifying \txNonce{}:}]
			% 	we impose $\txNonce_{i + \roffTxInitTxn} = \accNonce_{i + \roffTxInitAccDoingSenderPayForGas}$;
		\item[\underline{Justifying \txNonce{}:}]
			we impose
			\[
				\accNonce                              _{i + \roffTxInitAccDoingSenderPayForGas } =
				\txNonce                               _{i + \roffTxInitTxn                     } +
				\txNumberOfSuccessfulSenderDelegations _{i + \roffTxInitTxn                     }
			\]
		\item[\underline{Justifying $\txLeftoverGas$:}]
			cannot be set at the moment;
	\end{description}
	\iffalse
	I had to switch the order of phases from

	,   (TX_WARM) → (TX_AUTH) → (TX_SKIP ∨ TX_INIT) → ...
	,,,,,,,,,,,,,,,,,,,,,,,,,,,,,,,,,,,,,,,,,,,,,,,,,,,,,

	to 

	,   (TX_AUTH) → (TX_WARM) → (TX_SKIP ∨ TX_INIT) → ...
	,,,,,,,,,,,,,,,,,,,,,,,,,,,,,,,,,,,,,,,,,,,,,,,,,,,,,

	the reason for this switch is that it is easier to determine tx.requiresEvmExecution(), which decides whether we trace the TX_WARM phase, and whether we do TX_SKIP or TX_INIT, _after_ the TX_AUTH phase, given that the recipient account may have been delegated to a smart contract in the delegation list processing. Or the opposite.

	The real annoyance is that

	,   tx.requiresEvmExecution()
	,,,,,,,,,,,,,,,,,,,,,,,,,,,,,

	can no longer be determined from the transaction and the initial world state alone.
	\fi

	% \label{hub: initialization phase: common constraints: transactions supporting delegation lists must contain a TX AUTH phase}
	% We further impose that every \textbf{set code transaction} i.e. every \textbf{type 4 transaction}
	% begin with the processing of its delegation list.
	% \begin{description}
	% 	\item[\underline{Set code transactions must incorporate a \txAuth{}-phase:}]
	% 		\If $\txTypeSupportsDelegationLists _{i} = \true$ \Then
	% 		\[
	% 			\left\{ \begin{array}{lclr}
	% 				\txAuth          _{i - 1} & = & \true \\
	% 				\peekTransaction _{i - 1} & = & \true \\
	% 			\end{array} \right.
	% 		\]
	% \end{description}
	% \saNote{}\label{hub: initialization phase: common constraints: transactions supporting delegation lists must contain a TX AUTH phase}
	% Compare with
	% constraint~(\ref{hub: tx skip: user: transation processing: transactions supporting delegation lists start must contain a TX AUTH phase}).
	% Also recall
	% constraint~(\ref{hub: authorization phase: perspectives: the transation type must support delegation lists})
	% and the associated discussion.
